% =====================================================================
%  SAMPLE LAYOUT of the final report. It is NOT the report.
%  English edition of report_sample.tex: same structure, equations,
%  figures, tables and references; the notes are in English.
%
%  Real content (keep):   title, section headings, equations, figures,
%                         tables, reference list.
%  Yellow boxes (note):   a paragraph YOU write goes here. The box says
%                         what the paragraph must do. Delete the box
%                         when the paragraph is written.
%  Blue boxes (model):    example sentences about a DIFFERENT topic with
%                         invented numbers. Delete all of them.
%
%  The notes are instructions to the writer. They are not report text:
%  do not paste a note into the report and do not copy one with a few
%  words changed.
%
%  To see the bare layout without any box: change \notestrue to
%  \notesfalse on the line marked  <<< SWITCH  and compile again.
%
%  Compile:  tectonic report_sample_en.tex
%  Fonts: see the two cases in the comment of the font block below.
% =====================================================================
\documentclass{article}

\PassOptionsToPackage{numbers,sort&compress}{natbib}
\usepackage[preprint]{neurips_2024}
\makeatletter
\renewcommand{\@noticestring}{AA228/CS238 final project, Stanford University, Autumn 2026.}
\makeatother

\usepackage{amsmath,amssymb}
\usepackage{booktabs}
\usepackage{graphicx}
\usepackage{xcolor}
\usepackage{enumitem}
\usepackage{textcomp}
\usepackage{url}

% --- font block ---------------------------------------------------------
% This file has no Chinese text, so it does not load xeCJK.
% (a) tectonic or XeLaTeX (also on Overleaf): keep the two commands below
%     as they are. Without \setmainfont these engines silently set the
%     text in Latin Modern instead of Times, and the tables get wider.
% (b) pdfLaTeX: delete both commands; the style file then supplies Times
%     by itself.
\usepackage{fontspec}
\setmainfont{texgyretermes}[Extension=.otf,UprightFont=*-regular,BoldFont=*-bold,
  ItalicFont=*-italic,BoldItalicFont=*-bolditalic,Ligatures=TeX]
% --- end of the font block ------------------------------------------------

\usepackage{environ}
\usepackage[most]{tcolorbox}

\newif\ifnotes
\notestrue          % <<< SWITCH: \notestrue shows the boxes, \notesfalse hides them

\definecolor{noteback}{RGB}{255,249,219}
\definecolor{noteframe}{RGB}{150,105,0}
\definecolor{modelback}{RGB}{232,241,252}
\definecolor{modelframe}{RGB}{38,84,150}
\tcbset{sample/.style={breakable,enhanced,boxrule=0.4pt,arc=1.5pt,
  left=5pt,right=5pt,top=3pt,bottom=3pt,before skip=5pt,after skip=6pt,
  fontupper=\small}}

% \begin{note}[which paragraph, how long] ... \end{note}
\NewEnviron{note}[1][]{\ifnotes
  \begin{tcolorbox}[sample,colback=noteback,colframe=noteframe]
  {\bfseries\color{noteframe}You write\ifx\relax#1\relax\else: #1\fi}\par\smallskip
  \BODY
  \end{tcolorbox}\fi}

% \begin{model} ... \end{model}
\NewEnviron{model}{\ifnotes
  \begin{tcolorbox}[sample,colback=modelback,colframe=modelframe]
  {\bfseries\color{modelframe}Example sentences (a different topic: a traffic light; the numbers are invented; learn the style only)}\par\smallskip
  \BODY
  \end{tcolorbox}\fi}

\newenvironment{pts}{\begin{itemize}[nosep,leftmargin=1.15em,topsep=1pt]}{\end{itemize}}

% caption placeholder: the caption is yours to write
\newcommand{\capnote}[1]{\ifnotes{\color{noteframe}[Caption: you write] #1}\else Caption to be written.\fi}

% figures and tables are generated by code; show a frame if a file is missing
\newcommand{\fig}[2]{\IfFileExists{#1}{\includegraphics[width=#2]{#1}}%
  {\fbox{\parbox[c][1.5in][c]{0.9\linewidth}{\centering figure file not built yet:\\ \texttt{\detokenize{#1}}}}}}
\newcommand{\tab}[1]{\IfFileExists{#1}{\input{#1}}%
  {\fbox{\parbox[c][0.8in][c]{0.9\linewidth}{\centering table file not built yet:\\ \texttt{\detokenize{#1}}}}}}

\newcommand{\qprev}{q_{\mathrm{prev}}}
\newcommand{\rest}{\mathcal{J}}          % the set of forecasts not yet bought
\DeclareMathOperator{\E}{\mathbb{E}}

\title{2\textcent{} for Your Two Cents:\\ When Is an Unseen Forecast Worth Buying?}

% Use the names on the course roster and fill in the e-mail addresses.
\author{%
  Leonard Collomb\\ Stanford University\\ \texttt{(email)}
  \And
  Henry Chen\\ Stanford University\\ \texttt{(email)}}

\begin{document}

% =====================================================================
%  PAGE 0: how to read this sample (disappears with \notesfalse)
%  In this English edition the cover takes two pages; the paper's own
%  page numbers still start at 1 on the title page.
% =====================================================================
\ifnotes
\pagestyle{empty}
\begingroup\small
\begin{center}
{\Large\bfseries About this sample}\\[3pt]
{\small These two cover pages are not part of the paper. The paper starts on the page after them.\\
This is the English edition of \texttt{report\_sample.pdf} (the edition with Chinese notes). The content is the same.}
\end{center}

\subsection*{What this is}
\begin{pts}
\item A sample of the report with the layout already made. The layout, the sections, the equations, the figures, the tables and the reference list are real. The paragraphs of the text are empty: you write them. The course rule: the text of the report must be your own words. The conclusions must be your own judgment. So this sample does not contain any paragraph of report text written for you.
\item Write the English yourselves, directly. Do not give a note or a Chinese draft to translation software or to an AI to make the English. Text made that way is not your own words.
\item \textbf{The notes are instructions, not report text.} Write every sentence yourselves. Do not paste a note into the report. Do not copy a note with a few words changed.
\item A note sometimes gives one way to read a result (for example, ``the loss does not come from greedy''). That is only one possible reading. First check it yourselves against the figures and tables. Use it only if you agree.
\end{pts}

\subsection*{The three kinds of content}
\begin{note}[which paragraph, about how many words]
Yellow box: a paragraph that you write yourselves, in English, goes here. The box tells you the job of the paragraph, the numbers to use, and the textbook places to cite. When a paragraph is written, delete its whole box.
\end{note}
\begin{model}
Blue box: example sentences in English. They are about a different topic, and the numbers are invented. Look at how the sentences are built. Do not copy them, and do not reuse them with the nouns changed. Delete every blue box before you submit.
\end{model}
Content with no background color is real and can stay: the title, the section headings, the equations, the figures, the tables, the reference list.
\begin{pts}
\item The equations were written out exactly from the code. Derive each one by hand yourselves, until you can explain every term. You may change the notation.
\item The figures and tables are generated by code. Do not change a number by hand. You write the captions (the text that explains a figure or a table) yourselves.
\end{pts}

\subsection*{How to use it}
\begin{enumerate}[nosep,leftmargin=1.5em]
\item First read the whole sample once. See how long each section is and where the figures and tables go.
\item Order of writing: Sections 2, 3 and 4.1 $\rightarrow$ Sections 4.2 and 4.3 $\rightarrow$ Sections 5 and 6 $\rightarrow$ Section 1 $\rightarrow$ the abstract.
\item To see the layout alone, without the notes: in the source file, change \texttt{\textbackslash notestrue} to \texttt{\textbackslash notesfalse} and compile again. The reference list is then empty and BibTeX gives a warning, because all citations are inside the notes. This is normal.
\item When you have finished writing, send the text to Claude Code (the AI tool that made this sample; your teammate runs it). It checks the numbers, the logic and the textbook citations. It only points out problems. It does not rewrite your text.
\item Submission: a PDF. One team member submits it on Gradescope. The deadline is Friday, December 4, at 5 pm (Pacific time). Use the author names from the course roster. Fill in the e-mail addresses yourselves.
\item Before you submit, check: change \texttt{\textbackslash notestrue} to \texttt{\textbackslash notesfalse}; if you compile with pdfLaTeX, change the font block as its comment in the source file says; search the PDF for ``Caption to be written'' and for ``You write'' (neither may appear); fill in the e-mail addresses; count the pages.
\end{enumerate}

\subsection*{Words used in the notes}
These lines explain words to you. They are not text for the report: do not copy them into it.
\begin{pts}
\item \textbf{agent}: the program that makes the decisions each day. It does not train the forecast models. It only sees their outputs.
\item \textbf{free forecast} ($y_0$), \textbf{forecast A} ($y_A$), \textbf{forecast B} ($y_B$): three forecasts of $r$, the next-day price change of SPY. Free forecast: the agent sees it every day and pays nothing. Forecast A, forecast B: the agent sees one only after it buys it (0.1 bp each; each can be bought only once).
\item \textbf{position} ($q$): what the agent holds until the next day. \textbf{long} ($q=+1$): gains when the price rises. \textbf{flat} ($q=0$): holds nothing. \textbf{short} ($q=-1$): gains when the price falls.
\item \textbf{bp} (basis point): 1 bp $= 0.01\% = 10^{-4}$. All amounts are fractions of the money invested, reported in bp.
\item \textbf{net utility}: the score of one day. Equation (1) minus the prices paid for forecasts on that day. The three parts of Equation (1): the gain or loss, minus a risk penalty, minus the cost of changing position.
\item \textbf{belief}: the agent's probability distribution for $r$ (and for the forecasts not yet seen), given the forecasts already seen. \textbf{belief mean} ($m$): the mean of $r$ in that distribution.
\item \textbf{policy}: a rule that says which forecasts to buy each day. The five policies:
  \begin{itemize}[nosep,leftmargin=1.15em,topsep=0pt]
  \item \textbf{never buy}: never pays for forecast A or forecast B. It still chooses a position every day, by the same rule as the other policies, from the free forecast only.
  \item \textbf{always buy}: buys forecast A, then forecast B, every day.
  \item \textbf{random}: buys at random (the rule is in the third note of Section 4.1).
  \item \textbf{greedy}: looks one purchase ahead. It buys the forecast with the largest net value (Equation (4)) if that value is positive; if not, it stops.
  \item \textbf{exact}: plans both purchase decisions of the day together (Equation (5)).
  \end{itemize}
\item \textbf{always long}, \textbf{always flat}: two reference rows in Table 1. Position long (or flat) on every day; nothing bought.
\item $\boldsymbol{\lambda}$: the weight of the risk penalty in Equation (1). Larger $\lambda$ = more fear of risk. The experiments use $\lambda=1$ and $\lambda=5$.
\item \textbf{interval}: a 95\% block-bootstrap interval. A range around an average that shows how uncertain the average is. Made by resampling the days in blocks. If an interval contains 0, the notes say ``no conclusion''.
\item \textbf{fit period}: 1999-01-26 to 2014-12-31 (4{,}010 days); the belief is fitted on these days. \textbf{test period}: 2015-01-02 to 2026-11-12 (2{,}984 days); not used to fit the belief; the policies are compared on these days.
\item \textbf{predicted gain}: how much more the agent's model expects from following exact than from stopping immediately (no purchase). \textbf{realized gain}: the same difference as it really turned out. \textbf{same starting position}: on every day, both sides start from the position that never buy held the day before. Full definitions: fourth note of Section 4.1.
\end{pts}

\subsection*{Length (the course asks for 4--6 pages; ask the TA which parts count)}
\begin{center}\small\renewcommand{\arraystretch}{0.88}
\begin{tabular}{@{}lll@{}}
\toprule
Section & Words & What is in it\\
\midrule
Abstract & 150 & question, setting, method, result, conclusion\\
1 Introduction & 420 & two textbook sentences, two questions, earlier work\\
2 Problem Formulation & 220 & Figure 1, Equation (1)\\
3 Approach & 420 & Equations (2)--(5)\\
4.1 Setup and metrics & 360 & data, models, policies, two metrics\\
4.2 Greedy and exact & 230 & Table 1, Figure 2\\
4.3 Predicted and realized & 230 & Table 2, Table 3\\
5 Discussion and Limitations & 350 & limitations, future work\\
6 Conclusion & 70 & nothing new appears\\
Captions of the five figures and tables & about 200 & 30--50 words each\\
\bottomrule
\end{tabular}
\end{center}
\begin{pts}
\item Total: about 2{,}450 words, plus about 200 words of captions. The length was measured with placeholder text (filler text) of these word counts. Main text (to the end of the Conclusion): about 5.8 pages. With Contributions and the AI statement: about 6.2 pages. References: about 1 more page. In this test the AI statement had about 130 words and Contributions almost none. These page counts can be off by about $\pm 0.3$ page. Real text with many formulas will be longer.
\item So there is no extra space. Ask the TA (teaching assistant) this week: do the references, Contributions and the AI statement count in the 6 pages?
\item If Contributions and the AI statement also count: replace Table 3 by two sentences. This saves about 0.3 page; measured length then: about 5.9 pages. If the references count too: cut to about 1{,}700 words, remove Table 3, and write about 20 words per caption.
\end{pts}

\subsection*{How to regenerate the figures and tables}
Run these two commands in the \texttt{cs238/} folder:\\
{\footnotesize\texttt{python3 -m experiments.report\_figures --no-notes --out ../final\_report/figures}\\
\texttt{python3 -m experiments.report\_tables\ \ --no-notes --out ../final\_report/tables}}\\[2pt]
{\raggedright \texttt{--no-notes}: the tables come without the English note lines under them; you write the explanation in the caption. All result numbers are in \texttt{tables/numbers\_sheet.csv} (the numbers of the figures: \texttt{figures/figure\_numbers.csv}), each with its source. Copy the numbers in your text from there. The \texttt{meaning} column of the sheet is documentation of the code: copy the number, not the English wording of that column. The configuration values (model settings, dates) are in \texttt{cs238/configs/default.toml} and \texttt{cs238/results/tables/exp1\_lam1\_manifest.json}. Three kinds of numbers are not in these files; the notes that use them give the source: the download date of the data, the number of input features, the numbers of the October simulator.\par}
\endgroup
\clearpage
\pagestyle{plain}
\setcounter{page}{1}
\fi

% =====================================================================
%  THE PAPER
% =====================================================================
\maketitle
\suppressfloats[t]   % no figure above the title on the first page

\begin{abstract}
\begin{note}[abstract, about 150 words; write it after the rest of the paper is written]
Five sentences. Each sentence does one job only:
\begin{pts}
\item The question. Section 6.6 of the textbook has two sentences; you ask one question about each. First question: the textbook raises it itself. Second question: your own; the textbook sentence does not cover this case.
\item The setting: every day, before trading; the agent (the program that decides); two forecasts it has not seen yet, which it can pay for.
\item The method: greedy and exact, both computed; compared on the 2{,}984 days not used to fit the belief (the agent's probability model of $r$).
\item The result. Choose numbers from this list. Greedy and exact: same position on each of the 2{,}984 days. $\lambda=1$: predicted gain $+0.84$ bp per day, realized gain $-2.70$ bp (95\% interval $[-4.6,\,-1.0]$). $\lambda=5$: the interval contains 0, no conclusion.
\item The conclusion. What the two results mean together: judge this yourselves. Do not write a cause for the loss: the cause was not tested.
\end{pts}
\end{note}
\begin{model}
We ask whether a traffic signal that reacts to queue lengths leaves fewer cars waiting through a full cycle than a fixed timer does. In 30 simulated evenings per controller, 11.5\% of cars waited a full cycle under the reactive signal and 18.2\% under the timer; the spread across runs was about 2 points for each. We simulated a single intersection, so the result says nothing yet about a network.
\end{model}
\end{abstract}

% ---------------------------------------------------------------------
\section{Introduction}
% ---------------------------------------------------------------------

\begin{note}[paragraph 1, about 70 words: start from the textbook]
\begin{pts}
\item Describe the procedure that Section 6.6 of the textbook gives \citep{kochenderfer2022algorithms}. Which variable is observed at each step: the one with the largest value of information (the expected gain from seeing a variable before the decision). When the procedure stops: no further observation is worth making.
\item Two sentences of this section can be quoted directly, with quotation marks and page numbers:\\
page 121: ``The value of information is never negative.''\\
page 122: ``This greedy selection of observations is only a heuristic; it may not represent the truly optimal sequence of observations.''
\end{pts}
\end{note}

\begin{note}[paragraph 2, about 80 words: our problem]
\begin{pts}
\item Say what happens in one day, in one or two sentences. Point to Figure \ref{fig:day}(a).
\item Compare with an example that your classmates know: textbook Exercise 20.5 (page 424). In it: diagnostic tests one after another, then the decision to treat the patient or send the patient home. Matching parts: test = forecast; treat or not = position.
\item Give the differences from that example. Exercise 20.5: a test can be repeated. Here: each forecast can be bought only once. Exercise 20.5: tests independent of each other, given the patient's condition. Here: the two forecasts still correlated, even given $r$. You may also mention: the quantities here are continuous; one forecast is free.
\item Define each of these in one sentence when it first appears: SPY; position (long: gains when the price rises / flat: holds nothing, that is $q=0$ / short: gains when the price falls); bp.
\end{pts}
\end{note}

\begin{note}[paragraph 3, about 70 words: the two questions]
\begin{pts}
\item Question 1: greedy and exact, same choices or not?
\item Question 2: the agent's model predicts ``buying is better''. This gain on later data: did it appear or not?
\item Link the title to these two questions in one sentence. The title is a play on words, that is, words with a second meaning (here: paying money for another person's opinion). Price in the experiment: 0.1 bp per forecast.
\item State honestly: questions fixed after you saw the results; narrower than the proposal. Plan of the proposal: split the loss into three parts (planning, estimation, misspecification). Tested here: only the first part. Not separated: the other two parts.
\end{pts}
\end{note}

\begin{note}[paragraph 4, about 50 words: preview of the answers, and what you do not claim]
\begin{pts}
\item Question 1: greedy and exact, same position every day. Question 2, $\lambda=1$: the gain did not appear; the realized result was worse instead. Question 2, $\lambda=5$: no conclusion.
\item Say what ``never negative'' is about: expected utility under the agent's own model, before the price is subtracted. Say that you test a case outside this guarantee, and that you do not contradict the textbook. This sentence is required.
\item Do not explain the cause of the loss.
\item $r$, $q$ and $\lambda$ are not defined yet in this section (Section 2 defines them). In paragraph 2 and in this paragraph: either use words (next-day price change; the smaller and the larger weight of the risk penalty), or add a short pointer to Section 2.
\end{pts}
\end{note}

\paragraph{Related work.}\ifnotes\mbox{}\par\vspace{-3pt}\fi
\begin{note}[about 150 words: earlier work. The grading rubric gives 10 points for this paragraph and the reference list]
Four groups of papers. For each group: one or two sentences on what it did. Last sentence: where your work is different.
\begin{pts}
\item Value of information and ``one-step'' rules (rules that look only one step ahead): \citet{howard1966information} (origin of the concept); \citet{russell1991principles}; \citet{frazier2008kg} (a one-step rule for beliefs that are independent of each other, called knowledge gradient); \citet{frazier2010paradoxes} (marginal value of information not concave, so a one-step rule may stop too early); \citet{hay2012selecting} (definition of the myopic policy, myopic meaning short-sighted; remark that it tends to stop too early; Theorem 7: if the myopic policy decides to keep collecting, the optimal policy continues too).
\item Greedy compared with optimal. \citet{krause2009optimal}: direct comparison of the two in graphical models; proof that value of information is not submodular. \citet{golovin2011adaptive}: a greedy guarantee that needs adaptive submodularity, so it does not cover your objective function.
\item Sources or beliefs that are correlated with each other: \citet{frazier2009correlated}; \citet{clemen1985limits} (sources that depend on each other: the value of information falls quickly).
\item Combining forecasts, estimation error, disappointment after the decision: \citet{bates1969combination}; \citet{smith2006optimizer}; \citet{welch2008comprehensive} (note: it studies forecasts at monthly to five-year horizons, not daily ones).
\item Where you are different: do not say that ``comparing greedy and exact'' is new; \citet{krause2009optimal} and \citet{hay2012selecting} both compared them. What is different: the setting (buying forecasts for one trading position; Gaussian belief; a cost of changing position; each forecast can be bought only once) and Question 2 (predicted gain against realized gain on later data).
\item Cite only what you have read. Among the papers of this paragraph, five have no full text in the project folder: Howard 1966, Frazier and Powell 2010, Bates and Granger 1969, Clemen and Winkler 1985, Smith and Winkler 2006. Get each one through the university library before you cite it. Seven more cited entries also have no local full text (Hoerl and Kennard 1970, Breiman 2001, Friedman 2001, Pedregosa et al. 2011, Harris et al. 2020, Virtanen et al. 2020, Politis and Romano 1992); they are cited as the sources of methods and software. The page range of Politis and Romano 1992 comes from a secondary source: check it once before you submit.
\end{pts}
\end{note}

% ---------------------------------------------------------------------
\section{Problem Formulation}
% ---------------------------------------------------------------------

\begin{figure}[t]
\centering
\fig{figures/fig1_one_day.pdf}{\linewidth}
\caption{\capnote{One sentence on what the figure shows: (a) the decisions in one day, using the real day 2015-09-10, $\lambda=1$; (b) with yesterday's position long, the expected utility of the three positions as the belief mean changes. One sentence on what the reader should see: two forecasts bought; belief mean down from $+2.12$ to $-2.49$ bp; this is below the ``stay long'' threshold $+0.64$ bp; so the position changes to flat. You must state that you picked this day as an example: it is one of the 53 days with two forecasts bought; on a typical day only forecast A is bought (2{,}816 days). Also explain: open circle = belief before buying; filled circle = belief after buying; the same pair of points in both panels; curves in (b): cost of changing position already subtracted.}}
\label{fig:day}
\end{figure}

\begin{note}[paragraph 1, about 40 words: add the missing definitions. The timeline was told in the Introduction; do not repeat it here]
\begin{pts}
\item Define $r$, only once: the relative change of the dividend-adjusted closing price (adjusted close), from today's close to the close of the next trading day.
\item Give $r$ the name next-day price change. Reason: the ``return'' that the textbook defines (page 134) is the sum of rewards; another word avoids confusion.
\item Make clear that the agent is a program: it does not train the forecast models, it only sees their outputs.
\end{pts}
\end{note}

\begin{note}[paragraph 2, about 100 words: the formal statement of the problem. Choose one of the two versions below; do not mix them]
\begin{pts}
\item Version 1, POMDP (list the parts as the solution of Exercise 20.5 does). Hidden state: $r$ and the forecasts not yet seen. Observed part of the state: yesterday's position, which forecasts are already bought, the free forecast $y_0$. Actions: buy A, buy B, stop and choose a position. Observation: the value of the forecast that was bought. Reward: the price is paid at a purchase; Equation \eqref{eq:utility} is received at the stop. Add one last sentence: you solve it as a belief-state MDP, following Section 20.1 of the textbook.
\item Version 2, belief-state MDP (textbook Section 20.1, page 407). State: the belief mean and the belief variance, which forecasts are already bought, yesterday's position. Actions: the same as in Version 1. Transition: after a forecast is bought, belief update by Equation \eqref{eq:cond}. Reward: $-c_j$ at a purchase; at the stop, the expected utility of Equation \eqref{eq:stop} (textbook Equation (20.1): the expected reward under the belief).
\item Both versions must say: at most two purchase decisions per day, plus one choice of position; only link between the days: yesterday's position; no look-ahead across days by the planner. This last point is a limitation. Section 5 mentions it again.
\end{pts}
\end{note}

\begin{equation}
u(q, r) \;=\; q\,r \;-\; \lambda\,q^{2} r^{2} \;-\; \tau\,\lvert q - \qprev \rvert ,
\qquad q \in \{-1, 0, +1\}.
\label{eq:utility}
\end{equation}

\begin{note}[paragraph 3, about 80 words: explain Equation (1) term by term]
\begin{pts}
\item First term: the gain or loss. Second term: the risk penalty; larger $\lambda$ = more fear of risk. Third term: the cost of changing position, $\tau = 1$ bp.
\item Link to the textbook: first two terms = the quadratic utility of textbook Equation (6.3), $U(x)=\lambda x - x^{2}$, with argument $x = q\,r$ (the gain or loss of the day). Give the exact relation: $x - \lambda x^{2} = \lambda\,[(1/\lambda)\,x - x^{2}]$, that is, the textbook form with parameter $1/\lambda$, multiplied by $\lambda$. This factor cannot be dropped here: $\tau$ and the prices are in the same units. This step is your own derivation; the textbook does not write it this way. So use wording of the kind ``our $\lambda$''.
\item Define net utility: Equation \eqref{eq:utility} minus the prices paid for forecasts on that day. The tables and figures of Section 4 all use it.
\item Units: every quantity is a fraction of the capital (the money invested), reported in bp (1 bp $= 10^{-4}$).
\end{pts}
\end{note}

\begin{model}
The state is the pair of queue lengths, each between 0 and 20 cars. Every 5 s the controller takes one of two actions: keep the current green light or switch it. The reward is minus the number of waiting cars, so a larger reward means less waiting.
\end{model}

% ---------------------------------------------------------------------
\section{Approach}
% ---------------------------------------------------------------------

\subsection{Belief and its update}

\begin{note}[about 80 words]
\begin{pts}
\item Joint Gaussian of $(r, y_0, y_A, y_B)$: estimated once, on the first 4{,}010 days; not updated after that.
\item How it was estimated: means and variances = the sample values; correlations shrunk toward 0 (pulled part of the way to 0) with the Ledoit--Wolf shrinkage intensity (0.023) \citep{ledoit2004wellconditioned}. So write ``estimated''; do not write maximum likelihood. Explain the word shrinkage the first time you use it.
\item After each forecast that is seen: update with Equation \eqref{eq:cond}. In the textbook: Equations (3.25)--(3.26), Section 3.9 (page 63). Cite it.
\item One property worth writing: the variance after the update, no dependence on the value that was seen.
\item Before Equation \eqref{eq:cond}, first define the notation in one sentence: $o$ = the forecasts already seen; $\mu_r$, $\mu_o$ = the fitted means; $\Sigma_{rr}$ = the variance of $r$; $\Sigma_{oo}$ = the covariance of $o$; $\Sigma_{ro}$ = the covariance between the two (the textbook writes $A$, $B$, $C$; the letters are changed here because A and B already name the two forecasts).
\item Belief $b$ = the joint distribution of $r$ and the forecasts not yet seen, given $o$; its part for $r$: $\mathcal{N}(m, v)$.
\end{pts}
\end{note}

\begin{equation}
m \;=\; \mu_r + \Sigma_{ro} \Sigma_{oo}^{-1} (o - \mu_o), \qquad
v \;=\; \Sigma_{rr} - \Sigma_{ro} \Sigma_{oo}^{-1} \Sigma_{ro}^{\top}.
\label{eq:cond}
\end{equation}

\subsection{Choosing a position}

\begin{equation}
\E_b\!\left[u(q, r)\right] \;=\; q\,m - \lambda\,q^{2}\,(m^{2} + v) - \tau\,\lvert q - \qprev\rvert ,
\qquad
G(b) \;=\; \max_{q}\; \E_b\!\left[u(q, r)\right].
\label{eq:stop}
\end{equation}

\begin{note}[about 80 words]
\begin{pts}
\item In the textbook: the maximum expected utility principle, Section 6.4. $G(b)$ = the value of ``stopping now''.
\item Give both thresholds (a threshold is the level that $m$ must pass); which one applies depends on yesterday's position. Yesterday flat: go long only if $m - \lambda(m^{2}+v) > \tau$, which is roughly $m > \lambda v + \tau$. Yesterday long: stay long if $m$ is roughly above $\lambda v - \tau$.
\item Write out the arithmetic: standard deviation of $r$ about 128 bp, so $v \approx 1.64\times10^{-4}$ and $\lambda v \approx 1.64$ bp ($\lambda=1$). The two thresholds: 2.64 bp and 0.64 bp. For $\lambda=5$: about 9.2 bp and 7.2 bp.
\item Yesterday long or flat, and belief mean between the two thresholds (a band of width $2\tau$): the agent keeps yesterday's position. Figure \ref{fig:day}(b) shows the case $\lambda=1$, yesterday long.
\item Use these two thresholds to explain what $\lambda$ does. Section 5 uses them again.
\end{pts}
\end{note}

\subsection{Value of one more forecast}

\begin{equation}
\mathrm{VOI}_j(b) \;=\; \E_{y \sim b}\!\left[\,G\!\left(b^{j,y}\right)\right] - G(b),
\qquad
\mathrm{NV}_j(b) \;=\; \mathrm{VOI}_j(b) - c_j .
\label{eq:voi}
\end{equation}

\begin{note}[about 60 words]
\begin{pts}
\item First define the notation in one sentence: $j$ = a forecast not yet bought; $c_j$ = its price; $b^{j,y}$ = the belief after its value $y$ is seen; $\mathrm{NV}_j$ = the net value of buying it.
\item In the textbook: Equation (6.9) (page 121). Two changes: $EU^{*}$ replaced by $G$; the sum replaced by an expectation.
\item Link $\mathrm{VOI}_j(b) \ge 0$ to the textbook's ``never negative'': under the agent's own model, before the price is subtracted.
\item Gaussian model: the expectation in Equation \eqref{eq:voi} has a closed form (an exact formula; $G$ is a piecewise quadratic function of $m$); no sampling needed. One sentence is enough.
\end{pts}
\end{note}

\subsection{Greedy rule and exact solution}

\begin{note}[about 50 words: the greedy rule]
\begin{pts}
\item Describe the rule. For each forecast not yet bought: compute the net value $\mathrm{NV}_j$ of Equation \eqref{eq:voi}. Largest net value positive: buy that forecast. Otherwise: stop and choose a position.
\item In the textbook: the procedure in the last paragraph of Section 6.6. Call it by the textbook's name (greedy) and point to Section 6.6. The literature calls it myopic; mentioning that once is enough.
\end{pts}
\end{note}

\begin{equation}
V(b, \rest) \;=\; \max\Bigl\{\, G(b),\;\; \max_{j \in \rest}\,\bigl[\, -c_j + \E_{y \sim b}\, V\bigl(b^{j,y},\, \rest \setminus \{j\}\bigr) \bigr] \Bigr\},
\qquad V(b, \varnothing) = G(b).
\label{eq:bellman}
\end{equation}

\begin{note}[about 80 words: the exact solution]
\begin{pts}
\item First define the notation: $\rest$ = the set of forecasts not yet bought.
\item Equation \eqref{eq:bellman} = forward search over beliefs: textbook Section 22.2, Equation (22.1) (page 454), with discount 1, the sum replaced by an integral, and ``stop'' as the terminal action.
\item Depth 2 = the whole day. Reason: after both forecasts are bought, the only decision left is ``stop''. So: $G$ at the leaves exact; search not stopped early (not truncated).
\item How it is computed. Last step: expectation by the closed form of Equation \eqref{eq:voi}. First step in Equation \eqref{eq:bellman}: expectation by deterministic numerical integration; no random sampling. The integration rule: a composite Gauss--Legendre rule. The standardized range $[-12, 12]$ is cut at the kinks of the integrand (the points where its slope changes) and into panels at most 1 standard deviation wide; 201 nodes on each panel; at least 4{,}824 nodes in total. The \texttt{n\_quad = 201} of the configuration file is the number of nodes per panel. So ``exact'' means: exact except for this numerical integration.
\item Greedy and exact: a difference is possible only at the first decision. Reason: after one forecast is bought, the test for the remaining forecast is the same test in both. Write this point; Section 5 uses it.
\item Say that it is the exact solution under the agent's own model. Do not write only ``optimal'': it is optimal only inside the agent's model.
\end{pts}
\end{note}

\begin{note}[about 70 words: why this method. The 10 rubric points for Appropriateness of Approach are in this paragraph]
\begin{pts}
\item Forecasts not visible before they are bought, so: a sequential problem with information-gathering actions. Textbook page 122: the optimal sequence of observations can be found with the sequential decision methods of later chapters. That is exactly what you use.
\item Only two forecasts: so the exact solution can be computed, and only then is a direct comparison with greedy possible.
\item Gaussian model: closed form for the update and for the one-step value of information. So: no sampling error in the difference between greedy and exact.
\item Also give the reason for each choice in the evaluation, half a sentence each. Why the baseline (the policy the others are compared with) is never buy: same model, same rule for choosing the position, only no buying; so the difference comes only from the purchases. Why the comparison uses days not used for fitting: Question 2 asks about later data.
\end{pts}
\end{note}

% ---------------------------------------------------------------------
\section{Experiments}
% ---------------------------------------------------------------------

\subsection{Setup and metrics}

\begin{note}[paragraph 1, about 60 words: data]
\begin{pts}
\item Daily data of SPY; dividend-adjusted closing price. Source: Yahoo Finance, downloaded with the yfinance package; download date 2026-11-16 (UTC; source of this date: \texttt{cs238/data/raw/SPY\_daily\_manifest\_ci.json}). Data range: 1993-01-29 to 2026-11-13.
\item Fit period (the days used to fit the belief): 1999-01-26 to 2014-12-31, 4{,}010 days. Test period: 2015-01-02 to 2026-11-12, 2{,}984 days. Test period data: not used to fit the belief.
\item Explain the word held-out (data not used for fitting) the first time you use it.
\end{pts}
\end{note}

\begin{note}[paragraph 2, about 80 words: the three forecast models and the software. The rubric requires citations for algorithms and software that you did not develop yourselves]
\begin{pts}
\item Free forecast $y_0$: ridge regression \citep{hoerl1970ridge}. Forecast A: random forest \citep{breiman2001random}. Forecast B: gradient boosting \citep{friedman2001greedy}. All three with scikit-learn \citep{pedregosa2011scikit}.
\item Retrained every 21 trading days, with only the data available at that time; at least 1{,}260 days of training data.
\item Main hyperparameters (model settings chosen before training), half a sentence each (ridge: $\alpha=10$; random forest: 300 trees, depth 6; gradient boosting: 200 trees, depth 2, learning rate 0.05). One sentence on the input features: 16 of them (7 log returns over different windows, 3 realized volatilities, 6 others; source: \texttt{FEATURE\_COLUMNS} in \texttt{cs238/src/data\_pipeline.py}).
\item Other software: NumPy \citep{harris2020numpy} and SciPy \citep{virtanen2020scipy} for the numerical computation.
\item Give the location of the code in one sentence: repository address, commit, configuration file. The reader can find the other settings there. Before you cite it, commit what is not committed yet: 3 scripts, 4 result tables (CSV) and 3 test files. Also: the replay logs and the forecast archive are not in the repository (excluded by .gitignore), so a person who clones it cannot run the two scripts that generate the figures and tables directly. Say in one sentence how these files are obtained (for example: rebuilt by the GitHub Actions workflow of the repository; confirm this yourselves before you write it).
\end{pts}
\end{note}

\begin{note}[paragraph 3, about 80 words: the parameters and the policies compared]
\begin{pts}
\item Parameters: 0.1 bp per forecast; $\tau = 1$ bp; $\lambda \in \{1, 5\}$.
\item Price sweep: one extra run at each of 8 prices (0, 0.01, 0.03, 0.1, 0.3, 1, 3, 10 bp). Section 4.2 uses it, so introduce it here first.
\item Policies, half a sentence each: never buy (the baseline); always buy (forecast A first, then forecast B, every day); random (at each step it buys one more with probability 0.5, picked at random among those not yet bought; simulated only once); greedy; exact.
\item Fixed order (forecast A first, then forecast B; it decides only when to stop): choose one of two ways. Either do not mention it. Or mention it in one sentence: same purchases and same positions as exact on every day, so not in Table \ref{tab:policies}.
\item Two reference rows: always long, always flat.
\item State clearly: each policy enters the next day with its own position of yesterday.
\end{pts}
\end{note}

\begin{note}[paragraph 4, about 140 words: the two metrics. Define each one only once. Question 2 depends fully on Metric 2, so write it in full]
\begin{pts}
\item Metric 1: net utility of a policy minus net utility of never buy on the same day; averaged over days; unit bp/day; with a 95\% interval.
\item The interval: circular block bootstrap \citep{politis1992circular} (resampling whole blocks of days); blocks of 21 days; 2{,}000 resamples. The textbook does not cover it. Explain in one sentence why the resampling is by blocks (days next to each other are related).
\item Robustness check: 5 block lengths (5, 10, 21, 42, 63 days) $\times$ 20 bootstrap seeds, 100 intervals in total.
\item Report the interval endpoints with one decimal only, as in the tables. Reason: an endpoint changes with the random seed of the bootstrap, even in the first decimal, by a few tenths of a bp (example, exact at $\lambda=1$, 20 seeds: lower end from $-6.2$ to $-5.9$, upper end from $-1.1$ to $-0.9$). So the statement ``the interval excludes 0'' rests on the check with 100 intervals, not on the one decimal printed in the table.
\item Metric 2: define it with the notation of the paper. Predicted gain on one day: $V(b,\{A,B\}) - G(b)$, that is, how much more the agent's model expects from ``following exact'' than from ``stopping immediately''. Realized gain on that day: the net utility actually obtained by following exact, minus the net utility actually obtained by stopping immediately. Same starting position: both sides start every day from the position that never buy held the day before.
\item Say in one sentence why this is a test. The logic: if the fitted model is right, then mean realized gain = mean predicted gain (except for sampling error).
\item Say in one sentence why the numbers of Table \ref{tab:policies} and Table \ref{tab:gain} differ. Table \ref{tab:policies}: each policy on its own path of positions. Table \ref{tab:gain}: both sides start again every day from the same starting position.
\item Put the reading rule here: if an interval contains 0, say ``no conclusion''; do not say ``no effect''.
\end{pts}
\end{note}

\begin{model}
We count a car as delayed when it waits longer than one full signal cycle. For each controller we run 30 simulations and report the share of delayed cars, averaged over the runs, together with its standard deviation across runs.
\end{model}

\subsection{Greedy and exact purchasing}

\begin{table}[t]
\caption{\capnote{One sentence on what the table shows (the mean net utility of each policy, and its difference from never buy). One sentence on how to read the interval. One sentence on greedy and exact: same position on every day; the two rows differ only by the extra 0.1 bp paid on 1 day and on 8 days. The table itself has no notes, so all these explanations must be written here.}}
\label{tab:policies}
\centering\small
\tab{tables/tab1_policies.tex}
\end{table}

\begin{note}[paragraph 1, about 70 words: answer Question 1 directly]
\begin{pts}
\item Both values of $\lambda$: greedy and exact, same position on each of the 2{,}984 days.
\item Days with different purchases: 1 day for $\lambda=1$, 8 days for $\lambda=5$. All 9 days: greedy stops, exact buys forecast A; never the opposite case. Advantage that exact predicts on these days: below 0.001 bp each.
\item This direction: the same as in Theorem 7 of \citet{hay2012selecting} (their setting is different; write ``in a related setting'').
\item So treat the greedy row and the exact row of Table \ref{tab:policies} as one comparison, not as two independent pieces of evidence.
\end{pts}
\end{note}

\begin{note}[paragraph 2, about 80 words: read Table 1]
\begin{pts}
\item $\lambda=1$: exact vs never buy: $-3.36$ bp/day, interval $[-6.0,\,-1.0]$; always buy vs never buy: $-3.80$ bp/day, interval $[-6.6,\,-1.2]$.
\item random vs never buy: $-1.35$ bp/day, interval $[-2.7,\,0.0]$. A borderline result (just at the limit): the upper end of the interval is really $-0.005$, printed as 0.0 in the table; 74 of the 100 robustness intervals exclude 0; random itself simulated only once.
\item $\lambda=5$: the intervals of all the differences contain 0; no conclusion.
\item Two things you must point out. First: never buy and always long cannot be distinguished (4.42 vs 4.40, difference $-0.02$, interval $[-2.0,\,+1.9]$). Second: for $\lambda=5$, the highest point estimate (the single number, without its interval) is always flat.
\item The differences in the table are computed from unrounded numbers. So $3.06-4.42$ and the $-1.35$ of the table differ by 0.01: this is rounding, not an error.
\end{pts}
\end{note}

\begin{figure}[t]
\centering
\fig{figures/fig2_cumulative.pdf}{\linewidth}
\caption{\capnote{One sentence on what the vertical axis is (the cumulative difference in net utility from never buy, in bp). One sentence on what the reader should see (in which years the gap accumulates, that is, adds up; both panels use the same vertical axis). Explain that the cumulative value is a plain sum of the daily differences, not compounded. Steep drop at the end of the lines: real data; in the last 7 trading days, exact about 754 bp lower relative to never buy.}}
\label{fig:cumulative}
\end{figure}

\begin{note}[paragraph 3, about 80 words: greedy and exact also agree at other prices; when the gap accumulates]
\begin{pts}
\item Price sweep: 8 prices, two values of $\lambda$, 16 runs in total; greedy and exact, same position on every day.
\item Say that the reason is not ``the second forecast is rarely bought''. The numbers for this: $\lambda=1$, price 0: both forecasts bought on 89.7\% of days, positions still the same. (Price 0.1 bp: 1.8\%; 0.3 bp and above: 0.)
\item Figure \ref{fig:cumulative}: how the gap from never buy accumulates over time. Describe only what you see.
\item If you want to split the gap of Table \ref{tab:policies} into parts (sum of the 2{,}984 daily differences: $-10{,}018$ bp; mean $-3.36$ bp/day; the rounded $-3.36$ times 2{,}984 gives $-10{,}026$, and the gap of 8 bp is rounding): the 879 days with different positions: $-9{,}932$ bp. Among them, never buy long and exact flat: 333 days. Of these 333 days, exact already flat at the start of the day: 307. This point is optional.
\end{pts}
\end{note}

\begin{model}
Under the fixed timer 18.2\% of cars were delayed (standard deviation 2.1 points across 30 runs); under the reactive signal 11.5\% were (1.8 points). In heavy rain the two controllers differed by less than the run-to-run spread, so we do not rank them there.
\end{model}

\subsection{Predicted and realized gain}

\begin{table}[t]
\caption{\capnote{One sentence on what the columns are (predicted gain, realized gain, the difference of the two). One sentence on how the two sides are kept comparable (following exact, and stopping immediately: both sides start every day from the position that never buy held the day before). One sentence on what the last column is: for the difference ``realized minus predicted'', 100 intervals were computed (5 block lengths $\times$ 20 seeds); the column counts those that exclude 0.}}
\label{tab:gain}
\centering\small
\tab{tables/tab2_predicted_realized.tex}
\end{table}

\begin{note}[paragraph 1, about 70 words: answer Question 2 directly]
\begin{pts}
\item $\lambda=1$: predicted gain $+0.84$ bp/day; realized gain $-2.70$, interval $[-4.6,\,-1.0]$; difference of the two $-3.54$, interval $[-5.5,\,-1.8]$.
\item $\lambda=5$: predicted gain $+0.06$; realized gain $-0.32$, interval $[-1.5,\,+0.6]$; no conclusion.
\item Robustness, one sentence: $\lambda=1$: all 100 intervals exclude 0; $\lambda=5$: all 100 intervals contain 0.
\end{pts}
\end{note}

\begin{note}[paragraph 2, about 90 words: on which days the loss is made ($\lambda=1$). Describe only; do not give a cause]
\begin{pts}
\item This paragraph splits the $-2.70$ bp/day of Table \ref{tab:gain} into parts: $-8{,}065$ bp in total over the 2{,}984 days (the sum of the unrounded daily values; $-2.70 \times 2{,}984$ gives $-8{,}057$, the difference is rounding). Do not mix it with the $-10{,}018$ bp of Table \ref{tab:policies}.
\item Days on which the two sides hold different positions: only 374. Position differences on these days: $-7{,}770$ bp.
\item Prices paid for forecasts: only 295 bp in total. So point out where most of the loss is: not in the money paid for forecasts, but in the different position held after buying.
\item The side that follows exact: long on slightly fewer days, 71.2\% vs 73.5\%. This is the comparison of Table \ref{tab:gain} (same starting position). The long column of Table \ref{tab:policies} shows the opposite order (exact 77.88\%, never buy 73.53\%), because there each policy follows its own path. If you mention both, explain this.
\item The day of Figure \ref{fig:day}(a) can be used as an example, but say that it was picked. That day: two forecasts bought, then long changed to flat; next-day price change $+45.4$ bp; net utility of that day: never buy $+45.23$ bp, exact $-1.20$ bp.
\end{pts}
\end{note}

\begin{table}[t]
\caption{\capnote{One sentence on which two periods the two columns are. One sentence on what partial correlation means here (with the free forecast already known, how much a paid forecast is still related to $r$). One sentence saying that the numbers in parentheses are standard errors, and how they were computed (block bootstrap, blocks of 21 days).}}
\label{tab:periods}
\centering\small
\tab{tables/tab3_fit_vs_test.tex}
\end{table}

\begin{note}[paragraph 3, about 70 words: what each of the two periods looks like, see Table 3. Describe only; leave the explanation to Section 5]
\begin{pts}
\item Report every number together with its standard error. Partial correlation of forecast A: from $+0.030$ (0.025) to $-0.043$ (0.047). Of forecast B: from $+0.013$ (0.032) to $+0.104$ (0.059). Mean of $r$: from 2.75 (1.65) to 5.62 (1.73) bp.
\item Size of these three changes: only 1.2 to 1.4 standard errors. How computed: the change divided by $\sqrt{\mathrm{se}_{\mathrm{fit}}^{2}+\mathrm{se}_{\mathrm{test}}^{2}}$; this treats the two periods as independent. The other two rows of the table: standard deviation of $r$, 1.3 standard errors; correlation of $y_0$ with $r$, 1.9 standard errors, outside this range, so describe it separately. By your own reading rule, this is ``cannot be distinguished'', not ``changed''. Do not write that the predictive power of forecast A reversed.
\item Fit period already: the intervals of the two partial correlations contain 0.
\item Robustness: forecast A, in both periods: 0 of the 100 intervals exclude 0. Forecast B, in the test period: 27 of the 100, all of them at the two longest block lengths. So do not call it ``significant'' either. Free forecast $y_0$, in the test period: 100 of the 100 exclude 0.
\end{pts}
\end{note}

% ---------------------------------------------------------------------
\section{Discussion and Limitations}
% ---------------------------------------------------------------------

\begin{note}[paragraph 1, about 80 words: what the two answers mean together. Rubric: the conclusions must follow from the results]
\begin{pts}
\item The conclusion of this paragraph is for you to judge. Facts you can use:
\item Question 1: same positions at all prices. Limit of this fact: with only two forecasts, a difference between greedy and exact is possible only at the first decision. So it says little about ``is greedy good enough with more forecasts''.
\item Question 2, $\lambda=1$: the predicted gain did not appear. Question 2, $\lambda=5$: no conclusion. The thresholds can explain why: threshold to go long about 9.2 bp; never buy flat on 98.4\% of days; the side that follows exact in a different position from never buy on only 73 days (this is the comparison of Table \ref{tab:gain}; in the comparison of Table \ref{tab:policies}, each policy on its own path, the count is 81).
\item Meaning of the reference rows. The direct comparison is exact minus always long: at $\lambda=1$, $-3.34$ bp/day, interval $[-5.8,\,-1.0]$. So the loss here is also a loss against ``holding all the time''. (Never buy vs always long: the interval contains 0, no conclusion. Do not argue from ``cannot be distinguished''.)
\item One possible reading: what goes wrong here is not ``looking only one step ahead''. Use it only if you agree, and stop there; do not take the argument any further.
\end{pts}
\end{note}

\begin{note}[paragraph 2, about 80 words: what Table 3 and the thresholds can say together. Use ``is consistent with'', not ``is caused by'']
\begin{pts}
\item The relations that the agent's model relies on: already weak in the fit period, and not estimated precisely (Table \ref{tab:periods}).
\item At $\lambda=1$: fitted mean 2.75 bp, just above the ``enter'' threshold 2.64 bp, and about 2 bp above the ``stay long'' threshold 0.64 bp. Test period: belief mean before buying between these two thresholds on 38.6\% of days.
\item So: a small downward update, enough for an agent that is long to leave its position; to enter again, a higher mean needed. When you write this, say which explanation it is consistent with. Do not write it as a proven cause.
\item Caution, from Table \ref{tab:policies}: on its own path, exact is long on more days, not fewer (77.88\% vs 73.53\%). So do not write a general statement of the kind ``long less often after buying''. Where the loss is: days with never buy long and exact not long (never buy long and exact flat: 333 days, $-6{,}761$ bp in total).
\item \citet{smith2006optimizer}: about disappointment after a choice made by estimated values; directly related to Question 2. Read the paper first, then decide whether to cite it.
\end{pts}
\end{note}

\begin{note}[paragraph 3, about 70 words: the possible causes that were not separated. This paragraph decides whether the paper is honest]
\begin{pts}
\item Planning depth: tested; no difference here.
\item Three other possible causes, not separated: sampling error of the fitted model; change of the data between the two periods; wrong model family (Gaussian, linear).
\item Match them to the assumptions of the textbook. Sampling error and model family: the assumption that ``the model is given'' (start of Chapter 20, page 407: ``given models of the transitions, observations, and rewards''; the introduction of Part III says that this assumption is relaxed only there). Data that change over time: the stationary assumption of Section 7.1 (page 133).
\item The status update said that the cause was non-stationary skill. This claim was not tested (see the standard errors in Table \ref{tab:periods}). Correct it honestly in one sentence.
\end{pts}
\end{note}

\begin{note}[paragraph 4, about 70 words: limitations and assumptions. The text of the rubric requires them]
\begin{pts}
\item One asset; one split into periods.
\item Only two paid forecasts (so Question 1 tested only the first decision); prices artificial (set by you).
\item No look-ahead across days by the planner.
\item Belief: fitted only once; Gaussian.
\item Research questions: fixed after the results were seen.
\item One optional sentence; write it only if there is room and you have rerun it. October simulator: 14{,}175 settings; greedy worse than exact in 865 of them (6.10\%). Its parameters differ from the SPY experiment (in standardized units $\lambda s = 0.05$, $\tau/s = 0.05$, that is $\lambda=5$, $\tau=5$ bp when $s=1\%$; yesterday flat; belief mean from the trading boundary to one standard deviation above it). You must give these parameters with it. Do not put it next to the SPY numbers as a comparison. Source of these numbers: \texttt{analysis/a5\_output.txt}.
\end{pts}
\end{note}

\begin{note}[paragraph 5, about 50 words: future work. The text of the rubric requires it]
\begin{pts}
\item A paired control experiment: generate data from the fitted model, then estimate again and replay again. If predicted and realized agree there: a code error and some of the explanations can be excluded.
\item Put the uncertainty about the model parameters into the belief too. Sections 16.4--16.5 of the textbook have this idea, but they treat discrete transition probabilities, not Gaussian means and covariances. So cite them only as an idea.
\item Section 14.3 of the textbook (pages 288--289) is about the robustness of a policy to model error. You may cite it in one sentence. Do not expand on it: this chapter is not taught until December 3.
\end{pts}
\end{note}

\begin{model}
We tested one intersection with simulated arrivals. We did not test whether the reactive signal still helps when the arrival rate changes, so our result should not be read as a claim about real traffic.
\end{model}

% ---------------------------------------------------------------------
\section{Conclusion}
% ---------------------------------------------------------------------

\begin{note}[about 70 words, about five sentences]
\begin{pts}
\item One sentence repeats the questions. Two sentences give the two answers (one number each). One sentence says what this means for a reader of Section 6.6 of the textbook.
\item The last sentence names the main limitation and the next experiment. The rubric says ``The conclusion also includes discussion of limitations, assumptions, and future work'', so the conclusion must have this sentence.
\item Nothing new that did not appear before; no new numbers.
\end{pts}
\end{note}

% ---------------------------------------------------------------------
\section*{Contributions}
% ---------------------------------------------------------------------

\begin{note}[course requirement: after the conclusion, write what each person did]
Two or three sentences of facts: what each person really did (ideas, derivations, reviewing and testing code, running experiments, which sections each one wrote). Keep it consistent with the AI statement below: for code written by the AI, write here who reviewed it, who tested it and who made the decisions. Do not present that code as written by a person.
\end{note}

% ---------------------------------------------------------------------
\section*{Use of AI Tools}
% ---------------------------------------------------------------------

\begin{note}[course requirement: state which AI tools you used and what each one was used for. It must be honest and complete]
Below is a list of facts. You write the sentences yourselves:
\begin{pts}
\item Tool: Claude Code (the Claude models of Anthropic). If you also used others (translation software, a grammar checker, other chatbots), list them too.
\item Code and experiments: writing and debugging the code and the tests; running the experiments on GitHub Actions.
\item Analysis: designing and running the added analyses (the two reference rows always long and always flat; the robustness check over block lengths and seeds; the split of the loss by position); checking the numbers.
\item Figures, tables and layout: designing and generating the figures and tables; picking the example day of Figure \ref{fig:day}; this outline with its notes (the Chinese edition and the English edition), the LaTeX layout and the typeset equations.
\item Research questions and interpretation. Done by the AI: proposing and comparing several ways to state the research questions; recommending the present two questions; also proposing readings of the results and the list of ``what not to claim''. Which of these you adopted, and how you changed them, is your judgment. State it honestly.
\item Literature: finding papers and verifying the reference entries. Also: explaining concepts, making plans.
\item The text of the report: written by you. After the writing: check of the numbers, the logic and the citations by the AI (pointing out problems only, no rewriting). State this item too.
\item The statement in the status update was ``we used Claude Code to help write code; the experiment design, analysis, and all else are our own work''. The real use was wider than this sentence says. This time, write it in full.
\end{pts}
\end{note}

% ---------------------------------------------------------------------
%  References: every cited entry in refs.bib was checked on 2026-11-22
%  (Crossref, jmlr.org or the local PDF). One open item: the page range
%  of Politis and Romano (1992) comes from a secondary source.
% ---------------------------------------------------------------------
\bibliographystyle{unsrtnat}
\bibliography{refs}

\end{document}
